\section{Introduction}
\label{sec:intro}
Systems built from large language models are increasingly \emph{compiled} rather than hand-written: a developer states a task and a metric, and an optimizer searches over instructions, demonstrations, output formats or the structure of a multi-call program \citep{khattab2023dspy,opsahl2024optimizing,zhou2023large,yang2024large,shi2024efficient}. The search needs labelled questions, and in practice these are frequently the benchmark's own questions: the optimizer picks the configuration that scores best on a set of benchmark questions, and that best score on the same questions is then reported as the system's performance and used to compare systems. This is selection followed by evaluation on the data that drove the selection. Even when each configuration is only mildly better or worse than the others, the maximum over many noisy estimates is biased upward---the winner's curse \citep{smith2006optimizer}---so the reported number need not describe performance on new questions, and decisions based on it (``tuning gained $x$ points'', ``system A beats B'') can be wrong.

We ask two questions. \textbf{(Q1)} For a system whose configuration is searched on a finite set of questions and then used, how far is the reported score from its true accuracy on new questions, what controls the gap, and does it change conclusions drawn from reported scores? \textbf{(Q2)} Can reporting be made more reliable, and how should a reporting method be tested?

\textbf{Approach.} The gap is a property of the whole pipeline (data $\to$ search $\to$ report), so we study it by repeating the pipeline in worlds where the truth is known. We cache the correctness of 40--80 prompt configurations of real LLMs (Qwen2.5-0.5B/1.5B-Instruct, SmolLM2-1.7B-Instruct) on 450--750 real MMLU-Pro questions \citep{wang2024mmlupro}; taking the pool as the population gives exact true accuracies, and thousands of independent search episodes (five optimizer types imitating DSPy/MIPRO/best-arm-identification search logic) then measure the gap, its drivers, its effect on claims, and the bias/RMSE/coverage of eight reporting protocols. A short theory (\cref{sec:setting}) explains the results, and synthetic worlds isolate the factors.

\textbf{Findings.} (i) With $n{=}100$ dev questions the reported accuracy is $3.9$--$5.3$ points above the true accuracy of the deployed configuration---13--33\% of the accuracy itself---and above that of the \emph{best} configuration in the search space in 74--90\% of episodes, although the configurations differ by only 1--8 points (\cref{sec:exp}). (ii) The gap shrinks only as $\approx n^{-0.6}$, grows with the search effort (much faster than the selected configuration improves), with the number of effectively distinct configurations, and falls when one configuration is clearly best; a Gaussian plug-in predicts it within 10\%. (iii) It distorts conclusions: significance tests of a gain over a strong baseline fire 3--4$\times$ the nominal rate, mid-ranked baselines show lifts of $4.3$--$7.0$ points where the true lift is $0.3$--$2.0$, and systems tuned with unequal budgets are mis-ranked in magnitude. (iv) No protocol dominates: sample splitting is unbiased for any optimizer but noisy and costly; a search-replay bootstrap and empirical-Bayes shrinkage are far more accurate when their assumptions hold; a max-$t$ bound is always valid but wide.

\textbf{Contributions.} (1) A formal account of the post-search score with tight bounds and a plug-in predictor of its optimism (\cref{prop:opt,prop:cmp}). (2) A reusable meta-evaluation harness over cached real-LLM outcome matrices that measures bias, RMSE and coverage of reporting protocols with known truth. (3) A systematic measurement of this gap for LLM prompt/pipeline search on a real benchmark with real model outcomes, including its effect on lift and comparison claims. (4) A comparison of eight protocols and practical recommendations. Code, outcome matrices and all scripts are released.

\textbf{Conflict of interest disclosure.} None.
